\subsection{Sprint 0}

A proposta do projeto foi apresentada na reunião inicial com a \textit{stakeholder} Dayanne Silveira Antunes, marcando o início do projeto. Neste primeiro momento os AGES IV do time nos dividiram baseando-se nas preferências de cada aluno, seja pelo \textit{backend} ou \textit{frontend}. 

Sendo designado para o time do \textit{backend},  minha principal contribuição foi focada na modelagem dos esquemas conceituais e lógicos do banco de dados da aplicação. Essa tarefa acabou levando mais tempo do que o previsto, pois nossa visão das tabelas e das necessidades de armazenamento foi se ajustando ao longo da sprint.

Também me dediquei ao estudo das tecnologias NestJS e PrismaORM, já que nunca havia trabalhado com elas anteriormente. Para isso, foquei na construção de projetos pessoais e no uso dos materiais fornecidos pelo time.
\\

\fbox{%
    \parbox{0.95\linewidth}{%
        \setlength{\parindent}{15pt}
            \itshape Traçando um paralelo com minha primeira experiência na AGES, procurei tomar a iniciativa na comunicação com meus colegas para evitar possíveis problemas desde o início. Nosso time era composto por apenas quatro AGES II, em comparação com nove AGES I, o que me fez perceber rapidamente a necessidade de colaborar e buscar ajuda o quanto antes.

            Apesar da minha preocupação inicial com essa disparidade de experiência, o time se mostrou sempre unido, disposto e proativo. Os iniciantes se comunicavam bem e demonstravam grande vontade de aprender na AGES.
            
            Além disso, os líderes do time, tanto os AGES IV quanto nosso AGES III, prontamente se colocaram à disposição para ajudar com as tecnologias do projeto. 
    }%
}

